\begin{textsample}{POS Dim 1 – rn – Score 27.00 – rn/rn\_ef\_58.txt}  \label{ex:f1_pos_238}
ENSINO RELIGIOSO Ensino Fundamental – Anos Iniciais e Anos Finais Com o respaldo dos dispositivos legais citados, o Ensino Religioso apresenta -se como parte integrante do currículo escolar.
Ele tem como objeto de estudo o conhecimento religioso construído e acumulado ao longo da história, com ênfase na pluralidade e na diversidade das culturas e tradições religiosas.
A religião como fato \textbf{antropológico} e social \textbf{permeia} a vida dos cidadãos em qualquer sociedade, de todas as culturas, estabelecendo relações com os valores, a educação, a cultura e a política dos Estados.
Com o advento da República no Brasil, houve a separação entre a Igreja e o Estado, e toda a educação do cidadão deveria estar orientada com base no princípio da laicidade, de forma a garantir o critério da diversidade religiosa.
O Estado brasileiro laico respeita a liberdade religiosa, garante a consciência religiosa e se propõe a oferecer \textbf{uma} educação de qualidade e, desse modo, não se pode excluir a dimensão religiosa de \textbf{uma} proposta de formação integral.
Compreende -se \textbf{que} o ser humano se constrói a partir de \textbf{um} conjunto de relações tecidas em determinado contexto histórico-social [... ].
Nesse processo, o \textbf{sujeito} se constrói \textbf{enquanto} ser de imanência ( dimensão concreta, biológica ) e de transcendência ( dimensão subjetiva, simbólica ). ( BRASIL, 2017b, p. 436 ) A proposta de Ensino Religioso elaborada para o estado do Rio Grande do Norte fundamenta -se epistemologicamente pelas Ciências da Religião, conferindo aos estudos sobre o conhecimento religioso \textbf{um} viés de \textbf{abordagem} científica, sem cair no cientificismo, mas impedindo qualquer proposta de proselitismo ou confessionalismo.
Tal proposta compreende o Ensino Religioso segundo \textbf{um} modelo laico, pluralista, fundado no conceito de educação integral e democrática, com \textbf{forte} intenção de impedir qualquer prática catequética nas escolas do Rio Grande do Norte.
Ela se apresenta integrada às características pedagógicas e legais adotadas pela Secretaria de Estado da Educação e Cultura, em estreita articulação com a Base Nacional Comum Curricular ( BNCC ).
Nesse sentido, proporciona ao estudante as competências para “ reconhecer e cuidar de si, do outro, da coletividade, da natureza, \textbf{enquanto} expressões do valor da vida ” ( BRASIL, 2017b, p. 435 ) na observação de si, da família, da comunidade e do mundo onde está inserido e encontra os referenciais por meio dos quais organiza suas relações múltiplas.
No componente curricular Ensino Religioso, entende -se \textbf{que} religião e conhecimento religioso são patrimônios da \textbf{humanidade}, \textbf{uma} vez \textbf{que} histórica, política e culturalmente constituíram -se na \textbf{inter-relação} dos aspectos mais intrínsecos da vida humana, como : cultura, sociedade, economia, ética, estética, \textbf{sexualidade}, afetividade e espiritualidade.
Busca -se, ainda, apoiar o estudante no processo de educação integral para apropriar -se do conhecimento religioso, bem como conhecer as organizações religiosas das múltiplas e diversas culturas e tradições religiosas, colocando -se sempre em relação a outros campos e áreas do conhecimento.
Com isso, promove -se o reconhecimento da diversidade religiosa criada nas raízes culturais brasileiras, de modo a superar o preconceito religioso.
Portanto, proporcionam -se aos estudantes do Ensino Fundamental do Rio Grande do Norte subsídios para \textbf{que} entendam como os grupos sociais constituem -se culturalmente e, nessa contínua produção e legitimação de sua história, em simultaneidade desenvolvam e expressem sua relação com o sagrado.
Compete ao professor de Ensino Religioso possibilitar aos estudantes o livre acesso ao conhecimento e à compreensão das estruturas sobre as quais as religiões se constituem, bem como das interferências \textbf{que} causam no ambiente e no cotidiano, de modo \textbf{que} os estudantes obtenham, de forma ampla e completa, conhecimento crítico acerca das diferentes culturas e tradições religiosas presentes na sociedade.
Concordando com o pensamento de Edgar Morin ( 2007 ), acredita -se \textbf{que} os estudantes devem ser \textbf{educados} para a observação complexa da realidade, a fim de \textbf{que} tenham acesso a \textbf{um} conhecimento amplo e complexo sobre ela.
Cabe, ainda, inserir o estudante no campo das pesquisas sobre o conhecimento religioso e das tecnologias digitais da informação e comunicação ( TIDIC ) para possibilitar a ele saber lidar com a informação cada vez mais disponível, atuar com discernimento e responsabilidade nos contextos das culturas digitais, aplicar conhecimentos para resolver problemas, ter autonomia para tomar decisões, ser proativo para identificar os dados de \textbf{uma} situação e buscar soluções, conviver e aprender com as diferenças e as diversidades. ( BRASIL, 2017b, p. 14 ) Nesse sentido, para o professor ministrar as aulas, \textbf{exige} -se \textbf{que} trabalhe com competências e desenvolvimento de habilidades, de modo \textbf{que} possibilite transformar o conteúdo apresentado em conhecimento qualitativo e proveitoso para as relações dos seres humanos entre si e com outros seres vivos \textbf{que} compõem seu ambiente.
A fim de assegurar essas competências, julga -se \textbf{imprescindível} \textbf{uma} \textbf{sólida} formação inicial e contínua capacitação dos professores \textbf{que} ministram as aulas de Ensino Religioso.
O professor de Ensino Religioso \textbf{necessita} receber formação apropriada ao exercício pedagógico desse componente curricular e apropriar -se dos conhecimentos específicos \textbf{que} compõem essa área do conhecimento.
Com essa formação adequada, estará habilitado para oferecer aos estudantes a possibilidade de adquirir conhecimentos específicos e fazer relações desses conhecimentos com a transversalidade e a \textbf{interdisciplinaridade}, aspectos indispensáveis para \textbf{que} se compreenda a pluralidade e as diversidades sociocultural e religiosa, além da prática do respeito ao \textbf{que} é diferente.
Desse modo, o professor de Ensino Religioso deve ter como característica a “ busca do conhecimento das manifestações religiosas, [... ] a consciência da complexidade da questão religiosa e a sensibilidade à pluralidade ” ( FONAPER, 1997, p. 28 ).
Tais qualidades são essenciais para a \textbf{docência} desse componente curricular.
Para responder a tais desafios, o Ensino Religioso deverá sempre adotar metodologicamente a plena relação e o diálogo com o pensamento científico, bem como estar fundamentado no rigor do conhecimento construído na dinâmica do diálogo e no reconhecimento da alteridade.
Ele deverá também garantir \textbf{um} processo de ensino e de aprendizagem \textbf{que} estimule a construção do conhecimento mediada pelo diálogo, aberto ao debate e às hipóteses \textbf{que} se apresentem divergentes, não impedindo a dúvida, mas possibilitando a \textbf{exposição} competente de conteúdos \textbf{sistematizados}.
É por \textbf{uma} postura pedagógica \textbf{que} a religião e o conhecimento religioso serão estudados, para \textbf{que} haja o entendimento sobre as relações humanas produzidas na pluralidade e na diversidade culturais e religiosas das sociedades \textbf{contemporâneas}, marcadas pela laicidade, além de seu contexto originário, definido como secularização.
O estudo do conhecimento religioso tem como ponto de partida a(s) Ciência(s) da(s) Religião(ões), \textbf{embasada} em diversas áreas do conhecimento, em especial nas ciências humanas e sociais.
Pedagogicamente estrutura -se a partir de unidades temáticas distribuídas entre os nove anos do Ensino Fundamental, com suas respectivas habilidades e objetos de conhecimento.
A proposta pedagógica para o componente curricular Ensino Religioso organiza -se, de modo a facilitar a consulta, na seguinte ordem : Objetivos ; Competências específicas de Ensino Religioso para o Ensino Fundamental ; e Quadros organizadores, de acordo com os anos de escolaridade. \textbf{Salienta} -se \textbf{que} as \textbf{problematizações} e sugestões didáticas apresentadas são apenas orientações ; cabe ao professor a decisão de utilizá -las ou não.
Portanto, elas poderão, inclusive, ser modificadas, adaptadas e ampliadas para \textbf{que} atendam à realidade de cada escola e de cada turma.
O professor poderá, também, criar suas \textbf{problematizações} e sugestões didáticas tomando por base a própria experiência e o \textbf{que} melhor se adeque à realidade e à especificidade dos estudantes.
Assim sendo, observa -se \textbf{que}, em se tratando da validação da aprendizagem, busca -se a compreensão da avaliação na seguinte descrição : procedimentos e instrumentos de avaliação como orientação para a prática atrelada ao dia a dia da sala de aula, e não apenas a momentos pontuais.
Após os quadros organizadores curriculares ( 1o ao 5o ano e 6o ao 9o ano ), são apresentados os quadros de procedimentos e instrumentos de avaliação.
São sugeridos procedimentos e instrumentos de avaliação com base em como todos os estudantes \textbf{desempenham} certas \textbf{tarefas}, seja para diagnosticar a necessidade de realizar novas atividades para potencializar a aprendizagem, seja para dimensionar resultados ao final de \textbf{um} período/ciclo/etapa/bimestre.
Objetivos Considerando os marcos normativos e, em conformidade com as competências gerais estabelecidas no âmbito da BNCC, o Ensino Religioso deve atender aos seguintes objetivos : a ) proporcionar a aprendizagem dos conhecimentos religiosos, culturais e estéticos a partir das manifestações religiosas percebidas na realidade dos educandos ; b ) propiciar conhecimentos sobre o direito à liberdade de consciência e de crença, no constante propósito de promoção dos direitos humanos ; c ) desenvolver competências e habilidades \textbf{que} contribuam para o diálogo entre perspectivas religiosas e seculares de vida, exercitando o respeito à liberdade de concepções e ao pluralismo de ideias, de acordo com a Constituição Federal ; d ) contribuir para \textbf{que} os educandos construam seus sentidos pessoais de vida a partir de valores, de princípios éticos e da cidadania.

% matched lemmas: abordagem, antropológico, contemporâneo, desempenhar, docência, educar, embasar, enquanto, exigir, exposição, forte, humanidade, imprescindível, inter-relação, interdisciplinaridade, necessitar, permear, problematização, que, salientar, sexualidade, sistematizar, sujeito, sólido, tarefa, um
\end{textsample}
